% Additive recovery of N57, connected to the cylinder and signature
% proofs already included in the manuscript.
\subsection{The carry cocycle and the integer lift of the regional signature}
\label{carry27:seccion}

The regional signature retains the residue and quotient simultaneously.
The following construction makes the algebraic law of the latter
explicit. APP supplies integer evaluations on its two sheets;
TRIT retains orientation and regime; TPK transport produces the
regional bands and updates their memory. In a column of these
bands, Euclidean division determines the carry. The construction
uses the representatives \(\{0,1,2\}\) of \(\mathbb F_3\);
the oriented TRIT representation separately retains its sign
and transport quotient.

Let \((p,e,f)\in\{0,1,2\}^3\) be the entries of a regional
column. In this subsection, \(e\) denotes the ternary entry
of the propagation region. Define
\begin{equation}
\begin{aligned}
q&=[p+e+f]_3,& c&=(p+e+f-q)/3,\\
a&=[p+e-f]_3,& h&=(p+e-f-a)/3,
\end{aligned}
\label{carry27:division}
\end{equation}
where \([z]_3\in\{0,1,2\}\) is the remainder of any integer
\(z\). The variable \(a\) is the axial coordinate of this
column. The constant \(\alpha_{\rm HMT}\) retains its
subsequent dodecaphasic construction.

\begin{proposition}[Local composition of carries]
\label{carry27:cociclo}
For \(x,y,z\in\{0,1,2\}\), let
\[
\kappa_3(x,y)=\frac{x+y-[x+y]_3}{3},\qquad
\beta_3(x,y)=\frac{x-y-[x-y]_3}{3}.
\]
With \(r=[p+e]_3\), the two laws are
\[
c=\kappa_3(p,e)+\kappa_3(r,f),\qquad
h=\kappa_3(p,e)+\beta_3(r,f).
\]
Moreover,
\begin{equation}
\kappa_3(x,y)+\kappa_3([x+y]_3,z)
=\kappa_3(y,z)+\kappa_3(x,[y+z]_3).
\label{carry27:identidad}
\end{equation}
\end{proposition}

\begin{proof}
Adding \(p+e=r+3\kappa_3(p,e)\) to
\(r+f=q+3\kappa_3(r,f)\) gives the first formula.
Subtracting \(f\) and using
\(r-f=a+3\beta_3(r,f)\) gives the second.
Both sides of \eqref{carry27:identidad} equal
\((x+y+z-[x+y+z]_3)/3\). Composition therefore retains
the same quotient under either association of the sum.
\end{proof}

If \(s_3:\mathbb F_3\to\mathbb Z\) is the representative
section, then
\[
s_3(x)+s_3(y)-s_3(x+y)=3\kappa_3(x,y).
\]
This is the cocycle of the extension
\(0\to\mathbb Z\xrightarrow{\times3}\mathbb Z\to\mathbb F_3\to0\).
It arises from lifting a residue operation to integer
values. The visible coordinate and the memory of the
multiple of three are related by an exact identity.

\Needspace{10\baselineskip} % S27_LAYOUT_TABLE
\begin{longtable}{ccc|rrrr}
\caption{Complete law of the \(27\) regional columns.
The values of \(h\) are displayed as integers.}
\label{carry27:tabla}\\
\toprule
\(p\)&\(e\)&\(f\)&\(q\)&\(c\)&\(a\)&\(h\)\\
\midrule\endfirsthead
\toprule
\(p\)&\(e\)&\(f\)&\(q\)&\(c\)&\(a\)&\(h\)\\
\midrule\endhead
\bottomrule\endfoot
0&0&0&0&0&0&0\\
0&0&1&1&0&2&-1\\
0&0&2&2&0&1&-1\\
0&1&0&1&0&1&0\\
0&1&1&2&0&0&0\\
0&1&2&0&1&2&-1\\
0&2&0&2&0&2&0\\
0&2&1&0&1&1&0\\
0&2&2&1&1&0&0\\
1&0&0&1&0&1&0\\
1&0&1&2&0&0&0\\
1&0&2&0&1&2&-1\\
1&1&0&2&0&2&0\\
1&1&1&0&1&1&0\\
1&1&2&1&1&0&0\\
1&2&0&0&1&0&1\\
1&2&1&1&1&2&0\\
1&2&2&2&1&1&0\\
2&0&0&2&0&2&0\\
2&0&1&0&1&1&0\\
2&0&2&1&1&0&0\\
2&1&0&0&1&0&1\\
2&1&1&1&1&2&0\\
2&1&2&2&1&1&0\\
2&2&0&1&1&1&1\\
2&2&1&2&1&0&1\\
2&2&2&0&2&2&0\\
\end{longtable}

\begin{theorem}[The additional algebraic component of carry]
\label{carry27:rango}
Over \(\mathbb F_3\), the functions \(q,a\) are linear in
\((p,e,f)\). The reductions of \(c,h\) have polynomial
representations
\begin{align}
\bar c={}&2ef+2ef^2+2e^2f+2pf+2pf^2\notag\\
 &+2pe+pef+2pe^2+2p^2f+2p^2e,
\label{carry27:polinomio-c}\\
\bar h={}&2f^2+ef+2ef^2+e^2f+pf+2pf^2\notag\\
 &+2pe+2pef+2pe^2+p^2f+2p^2e.
\label{carry27:polinomio-h}
\end{align}
If \(V\) is the evaluation matrix of \(1,p,e,f\) on the
\(27\) columns, then
\[
\operatorname{rank}_{\mathbb F_3}V=4,\qquad
\operatorname{rank}_{\mathbb F_3}(V\mid\bar c)=5,\qquad
\operatorname{rank}_{\mathbb F_3}(V\mid\bar h)=5.
\]
\end{theorem}

\begin{proof}
Substituting the \(27\) triples from the table into the two
polynomials reproduces \(\bar c,\bar h\). The representation
of degree at most two in each variable is unique: a
polynomial of degree at most two in one variable that vanishes
at all three elements of \(\mathbb F_3\) is zero.
Applying this argument successively to the three variables
proves independence of the \(27\) monomials.

Evaluation at \(0,(1,0,0),(0,1,0),(0,0,1)\) proves
independence of \(1,p,e,f\). The function \(\bar c\)
vanishes at these four inputs and equals one at
\((1,1,1)\); it therefore lies outside their affine
evaluation space. For \(\bar h\), the first three
inputs force the constant, \(p\), and \(e\) coefficients
to zero, while \(\bar h(0,0,1)=2\) forces the
\(f\) coefficient to two. That affine function would
equal one at \((0,0,2)\), whereas
\(\bar h(0,0,2)=2\). Each additional column therefore
increases the rank by one.
\end{proof}

The proof specifies the information supplied by the
quotient. Residual superposition produces \(q,a\);
integer evaluation also incorporates \(c,h\).
For \(c\in\{0,1,2\}\), the residue determines its integer
value; for \(h\in\{-1,0,1\}\), it does so together with
this declared choice of representatives. The polynomials
describe the reductions; complete memory also retains
these lifts.

\subsubsection{Transport of the regional quotients}

Let \(M\) be an integer transport matrix already fixed
at a stage of the construction, with row states
\(x_t^\chi\) for \(\chi\in\{\pi,e,\varphi\}\).
Perform the following division coordinatewise:
\[
y_t^\chi=x_t^\chi M,\qquad
u_t^\chi=[y_t^\chi]_3,\qquad
x_{t+1}^\chi=(y_t^\chi-u_t^\chi)/3.
\]
Define \(X_t^+=x_t^\pi+x_t^e+x_t^\varphi\) and
\(X_t^\sigma=x_t^\pi+x_t^e-x_t^\varphi\).
Applying the preceding laws columnwise gives
\begin{equation}
\begin{aligned}
X_{t+1}^+&=\frac{X_t^+M-q_t}{3}-c_t,\\
X_{t+1}^\sigma&=\frac{X_t^\sigma M-a_t}{3}-h_t.
\end{aligned}
\label{carry27:transporte}
\end{equation}
Indeed, the sum of the three residues is \(q_t+3c_t\);
their signed combination is \(a_t+3h_t\). Substitution
of both identities into division of the states proves
\eqref{carry27:transporte}. The statement applies to
the declared integer transports. Regional emissions
and their continuity retain the coinductive readers
previously constructed.

As a complete evaluation on a six-column band,
\[
u^\pi=021020,\quad u^e=001220,\quad u^\varphi=100210
\]
give \(q=122120\), \(c=000110\), \(a=222000\), and
\(h=(-1,0,0,0,1,0)\). The orientation
\(\sigma=(1,1,-1)\) has doubled residue
\(2\sigma=(2,2,1)\) in \(\mathbb F_3^3\).
This last equality records the sign convention and
is distinguished from the six-column bands.

\subsubsection{Boundary memory and interval reading}

The local carry law couples to the previously established
cylinder refinement. One example exhibits the role of
the complete interval:
\[
\left[\frac1{729},\frac2{729}\right)
\cap
\left[\frac2{1000},\frac3{1000}\right)
=
\left[\frac2{1000},\frac2{729}\right)\ne\varnothing.
\]
By contrast, \(\lfloor1000/729\rfloor=1\).
The lower-endpoint reading returns block \(1\),
whereas interval compatibility also admits cell \(2\).
The boundary state retains both bounds needed to
decide that compatibility.

For one column, fixing \(q,a\) fixes
\[
f=2(q-a),\qquad p+e=q-f\quad\text{in }\mathbb F_3.
\]
Each of the nine fibres of \((q,a)\) contains exactly
three triples: choosing \(p\) determines \(e\) and
leaves \(f\) fixed. On six columns, before additional
constraints are imposed, the freedom is the product
of these six fibres and has cardinality \(3^6=729\).
Fixing the coordinate refers to alternatives at one
boundary. As depth increases, the register and its
memory change under transport. Local carry, cylinder
bounds, and regional incidence therefore act jointly
on the same arithmetic history.
